%%%%%%%%%%  scan idx 562  |  folio 549  %%%%%%%%%%
\SGAPrintPage{549}%

\sgaStableTarget{sga6.e10.thm.5-3-2}\textbf{Theorem 5.3.2.} \uline{The map} \(c_1\) (5.3.1) \uline{is an injective homomorphism, and the homomorphism}
\[
 \det:\Filt^1(X)\longrightarrow\Pic(X)
\]
\uline{induces, on passing to the quotient, a homomorphism that is a left inverse of} \(c_1\):
\[
 \alpha:\Gr^1(X)\longrightarrow\Pic(X).
\]
\uline{If every section} \(\underline{n}\) \uline{of the constant sheaf} \(\mathbb Z_X\) \uline{is bounded (for example, if} \(X\) \uline{is ``quasi-compact'', or if the final object} \(\underline{e}\) \uline{of} \(X\) \uline{is a finite sum of connected objects), then} \(c_1\) \uline{is an isomorphism and} \(\alpha\) \uline{is its inverse.}

\emph{Proof}: \(c_1\) is a homomorphism, since for every \(L,L'\in\Pic(X)\) one has the equation
\[
 (1-(LL')^{-1})-(1-L^{-1})-(1-{L'}^{-1})
 =
 -(1-L^{-1})(1-{L'}^{-1})\in\Filt^2(X).
\]

For the moment, admit that
\[
 \det(\Filt^2(X))=\{1\}\in\Pic(X)
\]
i.e., that \(\alpha\) is well defined. One then has
\[
 \alpha\circ c_1=\id_{\Pic(X)},
\]
since, by 5.2.2 b), one has
\[
 \det(1-L^{-1})=(\det(L-1))=L
\]
for every \(L\in\Pic(X)\). Next one shows that
\[
 c_1\circ\alpha=\id_{\Gr^1(X)}
\]
when every section \(\underline{n}\) of \(\mathbb Z_X\) is bounded. Since \(\Filt^1(X)\) is the abelian group generated by the elements \(\cl^{\bullet}(\underline{E})-\underline{n}\), where \(\underline{E}\) is a

%%%%%%%%%%  scan idx 563  |  folio 550  %%%%%%%%%%
\SGAPrintPage{550}%

locally free \(\underline{O}_X\)-module of rank \(\underline{n}\), it suffices, taking account of 5.2.2 b), to show that
\sgaStableTarget{sga6.e10.eq.5-3-3}\begin{equation}\tag{5.3.3}
 \det(\cl^{\bullet}(\underline{E}))-\cl^{\bullet}(\underline{E})+\underline{n}-1\in\Filt^2(X).
\end{equation}
Since \(\underline{n}\) is bounded, the final object \(\underline{e}\) of \(X\) decomposes into a \emph{finite} disjoint direct sum \(\coprod e_i\) of subobjects in such a way that \(\underline{n}\) is constant on each subtopos \(X/\underline{e}_i\) \((i=1,\ldots,p)\). Since the functor \(Y\rightsquigarrow K^{\bullet}(X/Y)\) \((Y\in\operatorname{Ob}(X))\) takes \emph{finite} direct sums to finite direct products, one may suppose that \(\underline n\) is constant on \(X\), say \(\underline n=n\in\mathbb Z\). Putting \(E=\cl^{\bullet}(\underline E)\), one finds (cf. Exposé V):
\[
 \det(\underline{E})=\lambda^n(E)=\gamma^n((E-n)+1)
 =
 \sum_{0\leq i\leq n}\gamma^i(E-n),
\]
This proves (5.3.3), whose left-hand side is equal to
\[
 \sum_{2\leq i\leq n}\gamma^i(E-n).
\]

It remains to prove the following

\sgaStableTarget{sga6.e10.lem.5-3-4}\textbf{Lemma 5.3.4.} \uline{For every ringed topos} \(X\) \uline{and every} \(x\in\Filt^2(X)\)\uline{, one has}
\[
 \det(x)=1.
\]

\emph{Proof}: Let \(\epsilon:K^{\bullet}(X)\to\Gamma(\mathbb Z)\) be the augmentation homomorphism. Then (6.2.8) and (6.5.3) below give, by linearity, the formula
\sgaStableTarget{sga6.e10.eq.5-3-5}\begin{equation}\tag{5.3.5}
 \det(xy)=(\det(x))^{\epsilon(y)}(\det(y))^{\epsilon(x)}
 \qquad(x,y\in K^{\bullet}(X)).
\end{equation}

%%%%%%%%%%  scan idx 564  |  folio 551  %%%%%%%%%%
\SGAPrintPage{551}%

It therefore suffices to show that, for every \(k>1\), one has
\[
 \det(\gamma^k(E-\underline{n}))=1,
\]
where, as above, \(E=\cl^{\bullet}(\underline E)\), \(\underline{n}=\epsilon(E)\). Now one has
\[
 \gamma^k(x)=\lambda^k(x+k-1)
\]
for every \(x\in K^{\bullet}(X)\), so that
\[
 \det(\gamma^k(E-\underline{n}))
 =
 \det(\lambda^k(E-\underline{n}+k-1))
 =
 \det\left(\sum_{p+q=k}\binom{-\underline{n}+k-1}{p}\lambda^q(E)\right).
\]
But (6.2.10) and (6.5.3) give
\[
 \det(\lambda^q(E))=(\det(E))^{\binom{\underline{n}-1}{\underline{q}-1}}
\]
for every \(q>0\), whence
\[
 \det(\gamma^k(E-\underline{n}))=(\det(E))^{\underline{r}},
\]
with
\[
 \underline{r}=\sum_{\substack{p+q=k\\ q\geq1}}
 \binom{-\underline{n}+k-1}{p}\binom{\underline{n}-1}{q-1}.
\]
Now \(\underline{r}\) is a section of \(\mathbb Z\), everywhere equal to the coefficient of \(t^{k-1}\) in \((1+t)^{k-2}\), and is therefore zero. The result follows.

\sgaStableTarget{sga6.e10.sec.6}\section*{6. Appendix: Calculation of Determinants of Locally Free Sheaves}

\sgaStableTarget{sga6.e10.sec.6.1}\subsection*{6.1.}
First consider an arbitrary ringed topos \(X=({C},\underline{O})\). For every positive section \(\underline{n}\in\Gamma(\mathbb Z_X)\) of the constant sheaf \(\mathbb Z\) and every \(\underline{O}\)-module \(M\in\Mod_{\underline{O}}(X)\), it is clear how to define the \(\underline{O}\)-modules

%%%%%%%%%%  scan idx 565  |  folio 552  %%%%%%%%%%
\SGAPrintPage{552}%

\sgaStableTarget{sga6.e10.eq.6-1-1}\begin{equation}\tag{6.1.1}
 M^{\underline{n}},\qquad \bigotimes^{\underline{n}} M,\qquad \bigwedge^{\underline{n}} M,\quad\text{etc.}
\end{equation}
The functors \((\ )^{\underline{n}}\), \(\bigotimes\limits^{\underline{n}}\) and \(\bigwedge\limits^{\underline{n}}\) take free (resp. locally free) modules to free (resp. locally free) modules.

More generally, consider sections \(\underline{n}_1,\ldots,n_r\) of \(\mathbb Z_X\) and a multifunctor (covariant, for definiteness)
\[
 T:\underline{\Mod}^{n_1}(\underline{O})\times\cdots\times
 \underline{\Mod}^{n_r}(\underline{O})\longrightarrow\underline{\Mod}(\underline{O})
\]
in \(r\) variables, where \(\Mod_{\underline{O}}^{\underline{n}}(X)\) denotes the category fibered over \(X\) of locally free \(\underline{O}\)-modules of rank \(\underline{n}\). Put
\[
 M=T(\underline{O}^{\underline{n}_1},\ldots,\underline{O}^{\underline{n}_r})
\]
and
\[
 G=\operatorname{G\ell}(\underline{n}_1)\times\cdots\times\operatorname{G\ell}(\underline{n}_r),
 \quad\text{where }\operatorname{G\ell}(\underline{n}_i)=\underline{\Aut}_{\underline{O}}(\underline{O}^{n_i}).
\]
It follows immediately that \(T\) defines a representation
\sgaStableTarget{sga6.e10.eq.6-1-2}\begin{equation}\tag{6.1.2}
 \alpha_T:G\longrightarrow\underline{\Aut}_{\underline{O}}(M).
\end{equation}
It is clear that, when \(M\) is itself locally free of rank \(\underline{m}\), for every composite functor
\[
 T'T,\ \text{where } T':\underline{\Mod}^{\underline{m}}(\underline{O})\longrightarrow\underline{\Mod}(\underline{O}),
\]
one has
\sgaStableTarget{sga6.e10.eq.6-1-3}\begin{equation}\tag{6.1.3}
 \alpha_{T'T}=\alpha_{T'}\,\alpha_T.
\end{equation}

Conversely, given the representation \(\alpha=\alpha_T:G\to\underline{\Aut}_{\underline{O}}(M)\), one recovers \(T(\underline{E}_1,\ldots,\underline{E}_r)\) (for every sequence

%%%%%%%%%%  scan idx 566  |  folio 553  %%%%%%%%%%
\SGAPrintPage{553}%

\(\underline{E}_1,\ldots,\underline{E}_r\) of locally free \(\underline{O}\)-modules of ranks \(\underline{n}_1,\ldots,\underline{n}_r\)) by gluing the free \(\underline{O}\)-module \(T(\underline{O}^{\underline{n}_1},\ldots,\underline{O}^{\underline{n}_r})\) on local trivializing charts by means of the representation \(\alpha\). More precisely, the sequence \(\underline{E}_1,\ldots,\underline{E}_r\) corresponds to a right \(G\)-torsor \(P=P_{\underline{E}_1,\ldots,\underline{E}_r}\) such that there is a canonical isomorphism
\sgaStableTarget{sga6.e10.eq.6-1-3.occ2}\begin{equation}\tag{6.1.3}
 \underline{E}_i\xrightarrow{\sim}P\overset{G}{\wedge}\underline{O}^{\underline{n}_i}\qquad(i=1,\ldots,r),
\end{equation}
where \(\overset{G}{\wedge}\) denotes the contracted product, and \(G\) acts on \(\underline{O}^{\underline{n}_i}\) through the \(i\)-th projection \(G\to\Gl(\underline{n}_i)\); then every representation
\[
 \alpha:G\longrightarrow\underline{\Aut}_{\underline{O}}(M)
\]
makes \(G\) act on \(M\) by \(\underline{O}\)-automorphisms, and thus defines a multifunctor \(T=T_\alpha\) such that
\sgaStableTarget{sga6.e10.eq.6-1-4}\begin{equation}\tag{6.1.4}
 T_\alpha(\underline{E}_1,\ldots,\underline{E}_r)=P_{\underline{E}_1,\ldots,\underline{E}_r}\overset{G}{\wedge} M.
\end{equation}
Moreover, for every composite representation \(\alpha=\alpha_1\circ\alpha_2\), one finds
\sgaStableTarget{sga6.e10.eq.6-1-5}\begin{equation}\tag{6.1.5}
 T_\alpha=T_{\alpha_1}\circ T_{\alpha_2}.
\end{equation}

\sgaStableTarget{sga6.e10.sec.6.2}\subsection*{6.2.}
For every locally free \(\underline{O}\)-module \(\underline{E}\) of rank \(\underline{n}\), we shall denote by \(\det(\underline{E})\) the invertible \(\underline{O}\)-module \(\bigwedge\limits^{\underline{n}} \underline{E}\). Keeping the notation of 6.1, the functor \(\det(T(\underline{E}_1,\ldots,\underline{E}_r))\) is determined, up to canonical isomorphism, by the corresponding representation
\sgaStableTarget{sga6.e10.eq.6-2-1}\begin{equation}\tag{6.2.1}
 \det\circ\alpha_T:G\longrightarrow{G}_m=\Gl(1).
\end{equation}

Suppose now that \(X\) is the ringed topos of sheaves on the ringed site \((\Sch/S)\) endowed with the étale topology. Then one can calculate \(\det(T(\underline{E}_1,\ldots,\underline{E}_r))\), or equivalently the representation \(\det\circ\alpha_T\), by means of the following
